\section{Related Work}
\label{sec:related_work}

\paragraph{Model collapse in iterative training.}

The phenomenon of iterative training collapse was formalized by \citet{shumailov2024curse}, who showed that recursively replacing real data with model-generated text causes progressive tail erosion across VAEs, GMMs, and language models, and by \citet{alemohammad2024mad}, who demonstrated quality--diversity tradeoffs in self-consuming generative loops for image models.
Theoretical analyses have since sharpened these results: \citet{dohmatob2024tale} recast collapse as a change in scaling-law exponents, showing that even small fractions of synthetic data alter effective scaling behavior, while \citet{dohmatob2025strong} proved that collapse can occur even when synthetic data constitutes a vanishingly small fraction of the training set.
The role of data accumulation has been a central theme.
\citet{dey2024breaking} proved that accumulating data across generations---rather than replacing it---prevents divergence of population risk in GMMs, HMMs, and language models.
\citet{kazdan2025collapse} systematically compared replace and accumulate regimes, finding that synthetic data value depends on the absolute quantity of co-present real data.
\citet{fu2025selfconsuming} provided generalization bounds for self-consuming loops, establishing that a constant fraction of real data suffices for convergence under accumulation.
On the calibration side, \citet{he2025golden} showed that optimal real-data weighting converges to the golden-ratio inverse ($1/\varphi \approx 0.618$) in Gaussian and GLM settings, and that the risk landscape is non-monotone in mixing ratio---a theoretical prediction with implications for regime-dependent canary behavior.
Beyond distributional statistics, \citet{guo2024linguistic} documented progressive decline in lexical, syntactic, and semantic diversity across iterative fine-tuning generations, providing the multi-level diversity metrics that serve as downstream targets in our lead-lag analysis.
\citet{schaeffer2025position} catalogued eight distinct and sometimes contradictory definitions of ``model collapse'' in the literature, underscoring the importance of explicit operational definitions---a concern our framework addresses by scoping canary prediction to diversity metrics specifically.

\paragraph{Canary metrics and collapse monitoring.}

Several recent works have proposed monitoring metrics that precede or accompany distributional degradation during iterative training.
\citet{du2026geometric} observe that calibration error rises before accuracy degrades in iteratively trained language models, interpreting this through a geometric-regulation lens.
\citet{chen2025scales} report a similar pattern across model scales and mixing ratios, noting that ECE increases of 75\% appear at 30\% synthetic contamination before accuracy drops become visible.
Beyond calibration, SIGMA~\cite{sigma2026spectral} introduces a dual-track spectral framework based on embedding Gram matrices, where divergence between deterministic and stochastic spectral bounds serves as an early collapse diagnostic.
\citet{kalinowski2026topology} proposes a topological monitoring approach using persistent homology features to detect multi-scale structural loss in representation space.
DCES~\cite{dces2026counter} uses center-edge sampling dynamics as a collapse signal and reports improvements in ECE and entropy preservation.
\citet{gambetta2024surplexity} introduces surplexity---a surprise-based metric that simultaneously detects collapse and filters training data---demonstrating that high-surplexity selection matches human-data baselines in maintaining distributional balance.

These works collectively demonstrate growing interest in collapse monitoring, yet they share methodological limitations that weaken their lead-lag claims.
None apply differencing tests to distinguish genuine temporal precedence from spurious co-trending; none cross-validate findings across multiple statistical methods; and none apply false discovery rate correction when testing multiple metric pairs.
Our work differs from the closest predecessor, \citet{du2026geometric}, in four respects:
(i)~we compare an entire family of canary metrics rather than a single metric pair;
(ii)~we cross-validate each lead-lag relationship across five independent statistical methods, quantifying inter-method sensitivity via Cohen's $\kappa$ and intra-method sensitivity via the Toda--Yamamoto augmentation order $d_{\max}$;
(iii)~we map detectability as a function of contamination regime rather than reporting a single aggregate result; and
(iv)~we treat method-selection sensitivity itself as a primary finding, showing that it is the dominant source of variation in existing lead-lag conclusions.

\paragraph{Early-warning signals and critical transitions.}

The theoretical foundation for canary-based prediction of system failure originates in ecology and climate science.
\citet{scheffer2009early} established that complex systems approaching tipping points exhibit generic early-warning signals---rising autocorrelation and variance due to critical slowing down---independent of system-specific mechanisms.
\citet{scheffer2012anticipating} extended this framework by combining empirical warning indicators with network-architecture analysis, arguing that structural and statistical signals together improve anticipatory power.
Crucially, \citet{boettiger2012limits} quantified the statistical detection limits of these signals, showing that early-warning power is often negligible in short time series---a constraint directly relevant to iterative training, where the number of generations (typically 5--20) is far smaller than typical ecological datasets.
\citet{truong2025entropy} connects entropy collapse to critical transitions, characterizing it as a first-order phase transition with universality classes that determine whether classical early-warning signals retain predictive power.
The model collapse literature has not yet imported these methodological tools: existing canary studies neither frame their claims in terms of detection power under finite samples nor test whether observed lead-lag patterns survive the short-series constraints that \citet{boettiger2012limits} identified as fundamental.

\paragraph{Lead-lag methodology.}

Granger causality~\cite{granger1969investigating} provides the standard econometric framework for testing whether one time series improves prediction of another beyond its own history.
\citet{shojaie2022granger} review classical and modern extensions, including VAR-based tests, kernel methods, and neural approaches for nonlinear dependencies \cite{marcinkevics2022neural}.
A key practical issue is non-stationarity: standard Granger tests applied to integrated (unit-root) series produce inflated rejection rates.
The Toda--Yamamoto procedure addresses this by augmenting the VAR lag order by the integration order and testing only the original lags, yielding valid asymptotic inference regardless of cointegration structure.
Despite wide use in economics and neuroscience, these methods are largely absent from the machine learning literature on iterative training.
Existing canary studies rely on visual inspection or naive cross-correlation, which cannot distinguish genuine temporal precedence from shared non-stationary trends.
Our framework applies five complementary methods---naive cross-correlation, differenced cross-correlation, threshold onset detection, permutation Granger causality, and Toda--Yamamoto Granger causality---and requires multi-method consensus before classifying a lead-lag relationship as robust, directly addressing the method-sensitivity problem that arises when any single test is taken in isolation.
Recent work on causal method reliability reinforces this multi-method rationale.
\citet{ruiz2026causalaudit} propose Causal-Audit, a pre-screening framework that validates causal assumptions before method selection and achieves 62\% false positive reduction; whereas their approach filters methods \emph{ex ante}, our consensus framework validates conclusions \emph{ex post} across methods already applied.
\citet{yi2026causalcompass} benchmark causal discovery robustness under eight misspecification scenarios and confirm that no single method dominates universally, directly supporting the multi-method rationale of our consensus framework.
Independently, \citet{sahoo2026calibration} show that calibration metrics are sensitive to training-distribution perturbation beyond the iterative collapse setting, suggesting that ECE-based canary signals may be confounded by non-collapse sources of calibration drift.
\citet{falahati2026curated} prove that pluralistic curation of synthetic data can prevent iterative collapse entirely, implying that canary detection is most relevant in uncontrolled data pipelines where such curation is absent.